\documentclass[10pt,a4paper]{amsart}
\usepackage[margin=19mm]{geometry}
\usepackage{amsmath,amssymb,mathtools}
\usepackage[scr=rsfs,cal=euler]{mathalfa}
\usepackage{xfrac}
\usepackage[unicode,hidelinks,bookmarks=false]{hyperref}
\newcommand{\ie}{i.e.\ }
\newcommand{\cf}{cf.\ }
\newcommand{\resp}{respectively\ }
\newcommand{\Subsection}{\subsection}
\providecommand{\nobreakdash}{\nobreak\hbox{-}}
\setlength{\emergencystretch}{2em}
\multlinegap0pt
\usepackage{bm}
\usepackage{bbm}
\usepackage{dsfont}
\usepackage{xspace}
\usepackage{cancel}
\usepackage{mathtools}
\usepackage{amsthm}
\makeatletter
\@ifundefined{lemma}{\newtheorem{lemma}{Lemma}}{}
\@ifundefined{proposition}{\newtheorem{proposition}[lemma]{Proposition}}{}
\makeatother
\def\ZZ{{\mathbb Z}}
\title[Critical set for surface diffeomorphisms revisited]{Critical set for surface diffeomorphisms revisited}
\author{Sylvain Crovisier, Enrique Pujals}
\date{Source: arXiv:2601.08208v2 (CC BY 4.0)}
\begin{document}
\maketitle

\noindent\emph{Source and licence.} Critical set for surface diffeomorphisms revisited. Version arXiv:2601.08208v2 (CC BY 4.0), distributed under the Creative Commons Attribution 4.0 International licence, as recorded in the arXiv licence metadata for this exact version and shown on the abstract page https://arxiv.org/abs/2601.08208v2. Full rights notice: licenses/arxiv-2601.08208-CC-BY-4.0.txt.\par\smallskip

\noindent\textit{Editorial scope.} Counted unit: the Proposition whose source label is \texttt{recurrent critical} in the article (source0.tex lines 1155--1162, source label \texttt{recurrent critical}), delivered together with the article's own complete original proof of it (source0.tex lines 1164--1205, one \texttt{proof} environment of 42 lines). No mathematics is written, repaired or re-derived: statement and proof are the article's own text line for line. Required context retained uncounted: none. Every definition, display and label of those lines is retained unchanged. Omitted from this package and not redistributed: 1--1154, 1163--1163, 1206--4481. Nothing inside a retained excerpt is omitted or altered: every retained body is delivered exactly as the article's lines are written, so no omission receipt of any kind accompanies this package. Citations: the retained excerpts carry no citation command at all, so no bibliography entry of the article is transcribed and no reference is redistributed. Reference adaptation: none was needed and none was made; every reference inside the delivered unit resolves inside it, so no unresolved reference or citation is produced. Printed slips kept literal: no printed word, symbol, hypothesis or number of the counted statement or of its proof was altered. Layout and class: the article's own class is \texttt{article} with packages including inputenc, amsfonts, amssymb, amsmath, amsthm, tikz, color, srcltx, epsfig, graphicx, geometry; the wrapper is the lane's pinned \texttt{amsart} editorial wrapper at 10pt on A4, which restates the theorem-like environments the retained text needs and adds the article's own macro definitions that the retained text needs (\texttt{\textbackslash ZZ}), with extra wrapper packages (bm, bbm, dsfont, xspace, cancel, mathtools). The delivered numbering is the unit's own. Assets: no retained span contains a figure environment, a graphics-inclusion command, a PDF-page inclusion, a table or a TikZ picture; no image file is copied into this package, the package declares no asset dependency and no image file is redistributed with it. Licence: version arXiv:2601.08208v2 of this article is distributed under the Creative Commons Attribution 4.0 International licence, as recorded in the arXiv licence metadata for this exact version and shown on the abstract page https://arxiv.org/abs/2601.08208v2. A full rights notice is delivered at licenses/arxiv-2601.08208-CC-BY-4.0.txt. Characters: every retained excerpt is pure ASCII, so the delivered engine, XeLaTeX with its default Latin Modern text font, reports no missing character.
\begin{proposition} \label{recurrent critical}

There exists a $C^1$-diffeomorphism  with an invariant compact set
conjugate to the shift on $\{0,1\}^\ZZ$ and with a single critical orbit.
The critical orbit can be taken  at any non-periodic point of $\{0,1\}^\ZZ$.

Moreover, the construction can be performed in such a way that the derivative of the projective cocyle does not grow exponentially. 
\end{proposition}

\begin{proof}
    Let us consider a dissipative diffeomorphism $f_0$ exhibiting a  classical horseshoe $K_0$ 
which is the maximal invariant set in an open domain $U$
and let us fix any non-periodic point $x\in K_0$.
We compose $f_0$ with a fibered rotation centered at $x$ with arbitrarily small
support and which almost sends $E^u_x$ to $E^s_x$.
We obtain a diffeomorphism $f_1$.
If the support $B_0$ of the rotation is small enough, a (narrow) cone field is still preserved,
the dynamics is still hyperbolic and conjugated to the shift,
the conjugacy is $C^0$-close to the identity,
the bundles $E^u,E^s$ are almost unchanged outside
the union of the $N_0$-iterates of the support $B_0$ of the rotation (where $N_0$ is large,
independent of the size of the support).
The new horseshoe $K_1$ is still the maximal invariant set in $U$.

We repeat the construction and get a sequence of diffeomorphisms $f_n$,
a sequence of horseshoes $K_n$, and a sequence of perturbation balls $B_n$.
%towers $T_nf_n^{-N_n}(B_n)\dots, B_n,\dots, f_n^{N_n}(B_n)$.
One can suppose that the diameter of $B_n$ and the angle of the rotation that we perform at step $n$
are less than $2^{-n}$ (after step $n$ the angle between $E^s$ and $E^u$
at $x$ is less than $2^{-n-1}$, so the rotation at step $n+1$ by an angle
smaller than $2^{-n-1}$ can still close the cone).
The convergence is thus strong enough, $f_n$ converges $C^1$ to
a diffeomorphism $f$ and $K_n$ to an invariant compact set $K$,
still the maximal invariant set in $U$,
and topologically conjugate to the shift.
Of course we still have dissipation.

The diffeomorphism $f_0$ preserved two transverse cone fields.
Outside these cones $g$ contracts (in the future or in the past depending on the cone).
Outside $B_0$ the cones are disjoint, so any point outside $B_0\cup f_0(B_0)$
does not satisfy the definition of critical set, for any diffeomorphisms $f_n$
(they all coincide there).
We want to repeat this argument so that after step $n$,
the critical set is contained in $B_n$. At the limit the critical set is reduced
to $x$ (which can be checked to be critical).
After step $n$, we consider two transverse invariant cone fields,  $C^s$ and $C^u$,
that are very thin, so that for $m_n$ large enough,
$g^{m_n}$ is expanding only inside $C^s$ and $g^{-m_n}$ is expanding only inside $C^u$.
Moreover after perturbation, we suppose that $f_{n+1}$ still sends $C^u$
disjoint from $C^s$, so that no point outside $B_{n+1}$ can be critical after step $n+1$.
\end{proof}

\end{document}
